An antenna's real performance is measured, not just simulated, and RF metrology
has its own methods and pitfalls.

\section{The Far-Field Condition}
Pattern and gain are defined in the far field, beyond
\[ R \ge \frac{2 D^2}{\lambda}, \]
where $D$ is the largest aperture dimension. Measurements are made on outdoor
ranges or in anechoic chambers whose absorber suppresses reflections.

\section{Gain Measurement}
Two standard methods: the \emph{gain-comparison} method references the antenna to
a calibrated gain standard, and the \emph{three-antenna} method solves for three
unknown gains from three pairwise transmission measurements --- needing no prior
standard.

\section{Pattern, Polarization, and Near-Field}
Rotating the antenna traces its radiation pattern and, with a polarized source,
its polarization and axial ratio. When a far-field range is impractical, a probe
scans the near field over a plane, cylinder, or sphere and a mathematical
transform computes the far-field pattern. Impedance and return loss are measured
directly on a VNA, and radiation efficiency by methods such as the Wheeler cap.
